\section{Ответы на контрольные вопросы}


\begin{enumerate}
	\item Оцените расстояние, на котором будет наблюдаться дифракция Фраунгофера на щели шириной 1 мм при длине волны 650 нм.
	
	\item Получите выражение для величины $I_{\max}$ в формуле \cref{eq:th:met:3} через интенсивность $I_0$ падающей на щель плоской волны.
	
	\item Выведите выражение для решеточного множителя в формуле \cref{eq:th:met:5}.
	
	\item Выведите формулу для распределения интенсивности $I(\theta)$ при освещении двух узких ($b \ll d$) щелей источником света конечного углового размера. Получите из нее соотношение \cref{eq:th:met:6}.
	
	\item Объясните наличие нулей в дифракционной картине на \cref{fig:2} с помощью векторной диаграммы.
	
	\item Сколько нулей и побочных максимумов находится между двумя соседними главными максимумами в дифракционной картине от решетки из $N$ щелей?
	
	\item Нарисуйте векторные диаграммы для первого и второго нулей картины при $N = 6$.
	
	\item С помощью векторной диаграммы найдите приближенно отношение высот первого побочного и центрального максимумов дифракционной картины при $N = 6$.
	
	\item В случае дифракции Фраунгофера на круглом отверстии как соотносится радиус отверстия с радиусом первой зоны Френеля?
\end{enumerate}